% Part: normal-modal-logic
% Chapter: syntax-and-semantics
% Section: relational-models

\documentclass[../../../include/open-logic-section]{subfiles}

\begin{document}

\olfileid{nml}{syn}{rel}

\olsection{اړيکيز مدلونه}

د نورمال مودالي منطقونو د معنٰی پوهنې بنسټيز مفهوم \emph{اړيکيز
مدل} دے۔ دا د نړۍو له يوه سټ څخه جوړ وي چې د دوه‌ځاييزې «لاسرسي
اړيکې» له لارې سره تړلې وي، او ورسره داسې يوه ټاکنه وي چې څرګندوي
کوم !!{propositional variable}s په کومو نړۍو کښې «رښتيا» ګڼل کېږي۔

\begin{defn}
  د بنسټيزې مودالي ژبې يو \emph{مدل} درې‌ګونې $\mModel{M}
  = \tuple{W, R, V}$ ده، چې په دې کښې
  \begin{enumerate}
  \item $W$ د «نړۍو» يو نا تش سټ دے،
  \item $R$ پر~$W$ د لاسرسي يوه دوه‌ځاييزه اړيکه ده، او
  \item $V$ داسې تابع ده چې هر !!{propositional
    variable}~$p$ ته د ممکنه نړۍو يو سټ $V(p)$ ټاکي۔
  \end{enumerate}
  که $Rww'$ وي، نو وايو چې $w'$ له~$w$ څخه \emph{ورته لاسرسے
    لري}۔ که $w \in V(p)$ وي، نو وايو چې $p$ په~$w$ کښې
  \emph{رښتيا دے}۔
\end{defn}

د اړيکيزې معنٰی پوهنې ستره ګټه دا ده چې مدلونه په ساده شکلونو
ښودل کېدے شي، لکه په \olref{fig:simple} کښې۔ نړۍوې د غوټو په بڼه
ښودل کېږي، او نړۍ~$w'$ هله او يوازې هله له~$w$ څخه ورته لاسرسے لري
چې له~$w$ څخه~$w'$ ته غشے وي۔ سربېره پر دې، که $w \in V(p)$ وي،
غوټه (نړۍ) په~$\mTrue{p}$ نښه کوو؛ که نۀ وي، په~\mFalse{p} ئې
نښه کوو۔ \olref{fig:simple} هغه مدل ښيي چې
$W = \{w_1, w_2, w_3\}$، $R = \{\tuple{w_1, w_2},\tuple{w_1,
  w_3}\}$، $V(p) = \{w_1, w_2\}$ او $V(q) = \{w_2\}$ لري۔

\begin{figure}
  \begin{center}
    \begin{tikzpicture}[modal]
      \node[world] (w1) [label={[align=right]right:\mTrue{p}\\ \mFalse{q}}]{$w_1$};
      \node[world] (w2) [label={[align=right]right:\mTrue{p}\\ \mTrue{q}},
        above right=of w1]{$w_2$};
      \node[world] (w3) [label={[align=right]right:\mFalse{p}\\ \mFalse{q}},
        below right=of w1] {$w_3$};
      \draw[->] (w1) to (w2);
      \draw[->] (w1) to (w3);
    \end{tikzpicture}
  \end{center}
  \caption{يو ساده مدل۔}
  \ollabel{fig:simple}
\end{figure}

\end{document}
